% GENERATED FILE. Edit canonical lessons, then run npm run generate:books.
% canonical-source-hash: fnv1a64:6ffe59c10954c1af
% canonical-lessons: HI-C39-dost, HI-C39-baccha, HI-C39-aadmi, HI-C39-aurat

\chapter{The People Around You}
\label{ch:hi-people-around-you}
\hlchaptermodality{39}

\begin{chapteropening}
\textbf{By the end of this chapter:} \emph{I can name the friend, the child, the man and the woman around me, I know that \hi{दोस्त} takes its possessive from the person while \hi{आँख} never does, and I can choose between \hi{औरत} and \hi{महिला} by register.}

Say the whole household with every possessive right — \hi{मेरा} \hi{पिता}, \hi{मेरी} \hi{माता}, \hi{मेरा} \hi{भाई}, \hi{मेरी} \hi{बहन}, \hi{मेरा} \hi{बच्चा}, \hi{मेरी} \hi{बच्ची} — then \hi{मेरा} \hi{दोस्त} or \hi{मेरी} \hi{दोस्त} depending on who the friend is, and say why Arabic ʿawra does not mean “woman” even though \hi{औरत} does.
\end{chapteropening}

\section[\hi{दोस्त}]{\hi{दोस्त} (dost) — the friend, who is \textquotedblleft{}the one chosen\textquotedblright{}}

\label{lesson:HI-C39-dost}



\begin{quote}
You can name a brother and a sister. Here is the person you picked rather than were given, and the word says exactly that.
\end{quote}

\subsection*{The letters in this word}
\textbf{\hi{दो}·\hi{स्त}} — an ending you have already assembled by hand.

\textbf{\hi{द}} with \textbf{\hi{ो}} gives \textbf{\hi{दो}}. Then \textbf{\hi{स्त}} — \textbf{\hi{स}} stripped of its built-in \emph{a} by the \textbf{virāma}, fusing with \textbf{\hi{त}} into one conjunct: the cluster inside \textbf{\hi{नमस्ते}}, the first conjunct you ever wrote.

So \textbf{\hi{मेरा} \hi{दोस्त}} is built entirely from strokes already yours.

\begin{grammarlens}[title={Grammar Lens: the noun that leans on who you mean}]
\textbf{\hi{दोस्त}} (\emph{dost}) = \textquotedblleft{}\textbf{friend},\textquotedblright{} and in the dictionary it is \textbf{masculine} — \textbf{\hi{मेरा} \hi{दोस्त}}.

But it bends. When the friend is a woman, ordinary modern Hindi says \textbf{\hi{मेरी} \hi{दोस्त}}. The word has not changed shape; the \textbf{possessive} followed the person. Both are current; neither is a mistake.

That is unlike everything so far. \textbf{\hi{आँख}} is feminine whatever you mean. \textbf{\hi{दोस्त}} takes its gender from \textbf{who the friend is}.
\end{grammarlens}

\begin{cousinweb}
\textbf{\hi{दोस्त}} is \textbf{Persian}, \textbf{\ar{دوست}} (\emph{dōst}); Old Persian had \textbf{dauštā}, an agent noun, \textquotedblleft{}\textbf{the one held dear}.\textquotedblright{}

Its root is \emph{\textbf{ǵews-}}, and the old meaning is not \textquotedblleft{}love\textquotedblright{} but \textquotedblleft{}\textbf{to taste, to try, to choose}.\textquotedblright{} Sanskrit inherited it as \textbf{\hi{जुष्}} (\emph{juṣ-}). Latin has \textbf{gustāre} — English \textbf{gusto}, \textbf{disgust} — and English \textbf{choose} is the same root.

So a \emph{dōst} is one you have \textbf{tried and chosen}; the narrowing to \textquotedblleft{}beloved\textquotedblright{} happened inside the Indo-Iranian branch alone.
\end{cousinweb}

\subsection*{You'll want to know — why it starts with a d}
The consonant at the front of this root came out as \textbf{z-} in Avestan and \textbf{d-} in \textbf{Old Persian}, by a regular sound law across the vocabulary.

Which is why \textquotedblleft{}friend\textquotedblright{} starts with \textbf{\hi{द}} in Persian and Hindi while its Sanskrit cousin \emph{juṣ-} starts with something else. An initial consonant is a passport stamp: it says which road the word came by.

\subsection*{Your turn}
Say these aloud:

\begin{itemize}
\raggedright
  \item \emph{merā dost} for a man, \emph{merī dost} for a woman
  \item the given and the chosen — \emph{merā bhāī, merī bahin, merā dost}
  \item \emph{namaste}, then \emph{dost} — one conjunct
  \item \emph{dōst}, \emph{juṣ-}, \emph{gustāre}, \emph{choose}
\end{itemize}

\begin{tcolorbox}[breakable,colback=teal!4,colframe=teal!35!black,title={Before you move on}]
How do you say \textquotedblleft{}my friend\textquotedblright{} about a woman? (\textbf{\hi{मेरी} \hi{दोस्त}} — the possessive follows the person.) What did the root behind \emph{dōst} first mean? (\textbf{To taste, try, choose}.) Why does Persian say \emph{d-} where Avestan says \emph{z-}? (\textbf{A regular sound law}.)
\end{tcolorbox}
\section[bacca]{\hi{बच्चा} (bacchā) — the child, who is a \textquotedblleft{}yearling\textquotedblright{}}

\label{lesson:HI-C39-baccha}



\begin{quote}
You have \textbf{\hi{पिता}} and \textbf{\hi{माता}}. Here is the third person at that table — and this word finally lets the masculine and feminine forms stand side by side.
\end{quote}

\begin{grammarlens}[title={Grammar Lens: a pair you can see}]
\textbf{\hi{बच्चा}} (\emph{bacchā}) = \textquotedblleft{}\textbf{child},\textquotedblright{} \textbf{masculine} — \textbf{\hi{मेरा} \hi{बच्चा}}.

Here the lean does real work at last. Swap the final \textbf{-\hi{ा}} for \textbf{-\hi{ी}} and you get \textbf{\hi{बच्ची}} (\emph{bacchī}), a girl child, \textbf{feminine} — \textbf{\hi{मेरी} \hi{बच्ची}}.

This is what \textbf{\hi{दरवाज़ा}} and \textbf{\hi{कुर्सी}} only hinted at: for \emph{this} class of noun \textbf{-ā}/\textbf{-ī} is not a lean but the machinery, and the word changes with the person. Which is exactly what \textbf{\hi{दोस्त}} would not do.

You already know how to ask this child's age: \textbf{\hi{कितने} \hi{साल} \hi{के} \hi{हो}?} — how many years someone \emph{is of}, rather than how many they \emph{have}.
\end{grammarlens}

\begin{cousinweb}
\textbf{\hi{बच्चा}} is \textbf{Persian}, \textbf{\ar{بچه}} (\emph{bacca}) — and the obvious story is wrong.

Persian \emph{bacca} did \textbf{not} come from Sanskrit \textbf{\hi{वत्स}} (\emph{vatsa}), \textquotedblleft{}calf, young one.\textquotedblright{} The two are \textbf{sisters}: one Indo-Iranian word inherited down two separate branches. Centuries later the Persian one was borrowed into Hindi — which had \emph{already} inherited the Sanskrit one as \textbf{\hi{बछड़ा}} (\emph{bachṛā}), \textquotedblleft{}calf.\textquotedblright{}

So Hindi holds both, the same ancient word by two roads — the pattern \textbf{\hi{दरवाज़ा}} and \textbf{\hi{द्वार}} showed you, running again, and the same Persian relay that brought \textbf{\hi{दोस्त}}.

Behind both stands \emph{\textbf{wet-}}, \textquotedblleft{}\textbf{year}.\textquotedblright{} Latin \textbf{vetus}, \textquotedblleft{}old,\textquotedblright{} is from it — English \textbf{veteran} — and so are Latin \textbf{vitulus}, \textquotedblleft{}calf,\textquotedblright{} and English \textbf{wether}, a yearling sheep.

A child, at the root, is a \textbf{yearling}. \textbf{\hi{पिता}} and \textbf{\hi{माता}} are named by relationship; \textbf{\hi{बच्चा}} is named by age.
\end{cousinweb}

\subsection*{Your turn}
Say these aloud:

\begin{itemize}
\raggedright
  \item \emph{bacchā} and \emph{bacchī}; then \emph{merā bacchā}, \emph{merī bacchī}
  \item the family — \emph{merā pitā, merī mātā, merā bacchā}
  \item the two roads — \emph{bacchā} from Persian, \emph{bachṛā} inherited; one word
  \item \emph{vetus}, \emph{vitulus}, \emph{wether} — and the meaning under all of them (year)
\end{itemize}

\begin{tcolorbox}[breakable,colback=teal!4,colframe=teal!35!black,title={Before you move on}]
What is the feminine of \textbf{\hi{बच्चा}}, and how does the possessive move? (\textbf{\hi{बच्ची}} — \emph{merā bacchā}, \emph{merī bacchī}.) Is Persian \emph{bacca} descended from Sanskrit \emph{vatsa}? (\textbf{No} — they are \textbf{sisters}; Hindi's inherited cousin is \textbf{\hi{बछड़ा}}.) What does the root mean? (\textbf{Year} — a child is a yearling.)
\end{tcolorbox}
\section[\hi{आदमी}]{\hi{आदमी} (ādmī) — the man, and the lean breaking backwards}

\label{lesson:HI-C39-aadmi}



\begin{quote}
The child has grown up. And this word does what none so far has — it ends in \textbf{-\hi{ी}} and is masculine anyway.
\end{quote}

\begin{grammarlens}[title={Grammar Lens: the lean breaks the other way}]
\textbf{\hi{आदमी}} (\emph{ādmī}) = \textquotedblleft{}\textbf{man, person},\textquotedblright{} and it is \textbf{masculine}. Say \textbf{\hi{एक} \hi{आदमी}} — it is also the ordinary word for \textquotedblleft{}person\textquotedblright{} in general, and stays masculine doing that. Ask \textbf{\hi{कितने} \hi{साल} \hi{के} \hi{हो}?} of him and you are asking how many years he \emph{is of}, the way Hindi counts age.

A warning before you reach for the possessive: \textbf{\hi{मेरा} \hi{आदमी}} is real Hindi but means what English's \emph{my man} means. Keep \textbf{\hi{एक} \hi{आदमी}}.

The gender still shows the instant anything agrees with the noun. Put it beside \textbf{\hi{कुर्सी}}: both end in long \textbf{-\hi{ी}}, one is feminine, one masculine, and no reasoning takes you from the ending to the answer. Beside \textbf{\hi{बच्चा}} / \textbf{\hi{बच्ची}}, where the ending \emph{does} carry the gender, the contrast is sharper.

So the lean has failed in \textbf{both} directions: \textbf{\hi{चाय}} and \textbf{\hi{किताब}} feminine with no \textbf{-\hi{ी}}, \textbf{\hi{आदमी}} masculine with one.
\end{grammarlens}

\begin{cousinweb}
\textbf{\hi{आदमी}} is \textbf{Arabic}, \textbf{\ar{آدمي}} (\emph{ādamī}), and it came the Persian way — the road of \textbf{\hi{कुर्सी}} and \textbf{\hi{किताब}}.

Now the ending. That \textbf{-ī} is not Hindi at all. It is the \textbf{Arabic ending meaning \textquotedblleft{}belonging to, of the kind of.\textquotedblright{}} Attach it to \textbf{\ar{آدم}} (\emph{Ādam}) and you get \textquotedblleft{}\textbf{of Adam}.\textquotedblright{} A human being, in this word, is one of Adam's.

\textbf{\ar{آدم}} is Hebrew \emph{\textbf{ʾādām}}, and here care is needed. You have probably heard that \emph{ʾādām} comes from \emph{\textbf{ʾadāmāh}}, \textquotedblleft{}ground\textquotedblright{} — man made of earth. Genesis sets the two side by side and means you to notice: a \textbf{deliberate pun}, and a famous one.

It is not the etymology. The prevailing view links the root \textbf{ʾ-d-m} to \textbf{redness} — behind Hebrew \emph{ʾādōm}, \textquotedblleft{}red,\textquotedblright{} and \textbf{Edom}. Beyond that, unsettled.

Worth keeping: \textbf{a wordplay a text makes on purpose is not where a word came from.}
\end{cousinweb}

\subsection*{Your turn}
Say these aloud:

\begin{itemize}
\raggedright
  \item \emph{ādmī}, then \emph{ek ādmī}
  \item \emph{kursī} feminine, \emph{ādmī} masculine; then \emph{merī kursī}, and stop
  \item the lean, and both ways it fails — \emph{chāy}, \emph{kitāb}, \emph{ādmī}
  \item what the \textbf{-ī} on \emph{ādamī} means
\end{itemize}

\begin{tcolorbox}[breakable,colback=teal!4,colframe=teal!35!black,title={Before you move on}]
\textbf{\hi{कुर्सी}} and \textbf{\hi{आदमी}} end alike — same gender? (\textbf{No}.) What is the \textbf{-ī} of \emph{ādamī} doing? (\textbf{Arabic \textquotedblleft{}belonging to\textquotedblright{}} — \textquotedblleft{}of Adam.\textquotedblright{}) Is \emph{ʾādām} from the word for ground? (\textbf{No} — Genesis is making a deliberate pun; the root is usually connected to \textbf{redness}.)
\end{tcolorbox}
\section[\hi{औरत}]{\hi{औरत} (aurat) — the word for \textquotedblleft{}woman,\textquotedblright{} and which one to use}

\label{lesson:HI-C39-aurat}



\begin{quote}
The last word of the chapter, and the one where Hindi hands you a choice rather than a word.
\end{quote}

\begin{grammarlens}[title={Grammar Lens: the pair, and the family complete}]
\textbf{\hi{औरत}} (\emph{aurat}) = \textquotedblleft{}\textbf{woman},\textquotedblright{} and it is \textbf{feminine}. Say \textbf{\hi{एक} \hi{औरत}} — and for the reason \textbf{\hi{एक} \hi{आदमी}} was the safe phrase, leave \textbf{\hi{मेरी} \hi{औरत}} alone: it means \emph{my woman}, a wife. \textbf{\hi{औ}} is one vowel letter, \emph{au}, standing alone.

Consonant ending, feminine — as \textbf{\hi{चाय}}, \textbf{\hi{किताब}} and \textbf{\hi{आँख}} were. Beside it \textbf{\hi{आदमी}}: \textbf{-\hi{ी}} ending, masculine. Neither adult's gender is reachable from spelling. Where you \emph{can} hear the agreement is on the child: \textbf{\hi{मेरा} \hi{बच्चा}}, \textbf{\hi{मेरी} \hi{बच्ची}}, the yearling who takes the ending seriously.

The family: \textbf{\hi{पिता}}, \textbf{\hi{माता}}, \textbf{\hi{भाई}}, \textbf{\hi{बहन}}, \textbf{\hi{बच्चा}}, \textbf{\hi{आदमी}}, \textbf{\hi{औरत}} — and the chosen one, \textbf{\hi{दोस्त}}.
\end{grammarlens}

\begin{cousinweb}
\textbf{\hi{औरत}} is \textbf{Arabic}, \textbf{\ar{عورة}} (\emph{ʿawra}), on the three-consonant root \textbf{ʿ-w-r} — built the way \textbf{k-t-b} built \textbf{\hi{किताब}} — whose senses run to \emph{a weak point, a gap in a defence, something left uncovered}.

Then the fact that changes what you do with all that: \textbf{Arabic ʿawra does not mean \textquotedblleft{}woman.\textquotedblright{}} That sense is a \textbf{Persian} development. Persian moved the noun to mean a woman or a wife and handed \emph{that} to Hindi and Urdu, on record from around 1420 — the same Persian relay that carried \textbf{\hi{दोस्त}}.

So the harsh root sense is not a claim Arabic makes about women. It is the earlier life of a word another language repurposed.
\end{cousinweb}

\subsection*{You'll want to know — three words, and which to say}
Three words, and the choice is \textbf{register}, exactly as for \textbf{\hi{किताब}} / \textbf{\hi{पुस्तक}} / \textbf{\hi{पोथा}}.

\begin{itemize}
\raggedright
  \item \textbf{\hi{औरत}} — Perso-Arabic, plain and everyday. Common in speech; some speakers avoid it formally, partly because of the origin above.
  \item \textbf{\hi{महिला}} (\emph{mahilā}) — a Sanskrit borrowing, the \textbf{respectful, public} word: signs, announcements, official prose. Its Sanskrit origin is \textbf{uncertain}; one serious proposal traces it to Dravidian.
  \item \textbf{\hi{स्त्री}} (\emph{strī}) — Sanskrit, older and literary. Secure back to Indo-Iranian, no further.
\end{itemize}

\textbf{\hi{महिला}} where you would write, \textbf{\hi{औरत}} where you would talk.

\subsection*{Your turn}
Say these aloud:

\begin{itemize}
\raggedright
  \item \emph{ek aurat}, then \emph{ek ādmī} — neither ending gives its gender away
  \item \emph{merā pitā, merī mātā, merā bhāī, merī bahin, merā bacchā, merī bacchī}
  \item \emph{merā dost}, then \emph{merī dost}
  \item every feminine noun you own — \emph{kursī, chāy, kitāb, āṁkh, bacchī}
\end{itemize}

\begin{tcolorbox}[breakable,colback=teal!4,colframe=teal!35!black,title={Before you move on}]
Does Arabic \textbf{\ar{عورة}} mean \textquotedblleft{}woman\textquotedblright{}? (\textbf{No} — that sense is \textbf{Persian}.) Which goes on a public sign, which in conversation? (\textbf{\hi{महिला}} written; \textbf{\hi{औरत}} spoken.) Your family, minding every possessive? (\emph{Merā pitā, merī mātā, merā bhāī, merī bahin, merā bacchā}.)
\end{tcolorbox}
